\subsection{Проблемы базовых органов реле направления мощности}
\label{sec:ch1_base_organs_problems}

В случае отсутствия двигательной нагрузки на секции шин, при возникновении внешних аварий ток практически сразу пропадает. Однако ввиду математических особенностей расчёта векторов тока и напряжения с применением Фурье преобразования -- обнуление сигналов используемых в расчёте РНМ происходит с запаздыванием равным периоду сети, то есть с выдержкой порядка 20~мс. Данную особенность можем видеть на рисунке 1, здесь фаза А представлена амплитудой первой гармоники, а фаза С -- мгновенными значениями.

\begin{figure}[htbp]
    \centering
    \includegraphics[width=0.85\textwidth]{image5.png}
    \caption[Поведение амплитуды тока при расчёте первой гармоники]{Поведение амплитуды тока при \commentnote{DISS-005-C0015}{Feel}{Понимаю, что фаза А, потому что желтая. Токи подписаны одинаково. Печать может быть ч/б, лучше подписать по-разному}{\commentnote{DISS-005-C0016}{Евдаков Алексей}{Да, сделаем, наверное, прост симулируем, тут не нужен прям реальный сигнал, лишь описание сути.\newline + сделать без колебаний -- ровную.}{расчёте}} первой гармоники.\newline{\small(Файл находится в «\_вспомогательные файлы\textbackslash{}Краснодар Водоканал»)}}
    \label{fig:ch1_img5}
\end{figure}

Аналогичное поведение может наблюдаться и на подстанциях с двигателями при отключении питающих линий. Причём в ряде случаев, \commentnote{DISS-005-C0017}{Feel}{Нужно пояснить. Сигналов токов в фазах В и С? А что произошло? ОЗЗ? Тогда токи вроде и не должны до нуля упасть, да и напряжения есть (треугольник сохранился)}{\commentnote{DISS-005-C0018}{Евдаков Алексей}{Да, совсем не то вставил.\newline Вот новая картинка, на которой видно, что ток не обнулился.\newline Тип аварии -- внешнее двухфазное замыкание в сети 110кВ (за трансом).}{обнуление}} сигнала вообще может не наблюдаться, что можно видеть из рисунка 2.

\begin{figure}[htbp]
    \centering
    \includegraphics[width=0.85\textwidth]{image6.png}
    \caption[Полный вид осциллограммы произошедшего события по одной из секций]{Полный вид осциллограммы произошедшего события по одной из секций\newline{\small(Файл находится в «\_вспомогательные файлы\textbackslash{}ДКС-1»)}}
    \label{fig:ch1_img6}
\end{figure}

При этом расхождение по углу не всегда будет достаточным для срабатывания БАВР. И в результате функционирование БАВР будет происходить либо с заметной задержкой, которая заранее не определена, либо вообще может не произойти.

Например,\commentnote{DISS-005-C0019}{Евдаков Алексей}{13000705}{\commentnote{DISS-005-C0020}{Евдаков Алексей}{Там же есть 13000704 -- там просадка была более долгая, но не такая сильная, и БАВР по напряжению мог бы отработать позже. В общем, можно будет подумать и в случае чего и её включить}{}} на рисунке 3 мы можем видеть, что РНМ отпустил только через 470~мс, что очень долго. Данное поведение стандартного РНМ обусловлено тем, что это несимметричное снижение напряжения с достаточно большим остаточным напряжением, и как следствие -- снижения потребления мощности практически не происходит (лишь кратковременно). Так как для работы БАВР в некоторых устройства требуется, чтобы условия сформированные -- приводили к работе в течении 200~мс, что не было выполнено в данном случае из-за РНМ -- то БАВР перешёл в логику работы штатного АВР (0,6$U_{\text{ном}}$ и 1~с), что тоже не успело сформироваться. А значит ни БАВР, ни АВР вообще не отработали, поэтому весь аварийный процесс для 1СШ продолжался 1,57~с, что очень и очень много. При этом, в данных случаях -- функционирование БАВР требуется.

Отметим, что подобная логика работы БАВР (с блокировкой через 200~мс) не является общепризнанной и не встречается на всех устройствах. И в данном случае можем видеть, что она привела к отказу \commentnote{DISS-005-C0021}{Евдаков Алексей}{Это не совсем РНМ касается, но на самом деле важное примечание. Поэтому пока жёлтым.}{функционирования}. Но это является отдельным вопросом, и поэтому подробно разбираться не будет, отметим лишь, что к этой проблеме исходно привела блокировка от РНМ.

\begin{figure}[htbp]
    \centering
    \includegraphics[width=0.85\textwidth]{image8.png}
    \caption[Полный вид осциллограммы произошедшего события]{Полный вид осциллограммы произошедшего события\newline{\small(Файл находится в «\_вспомогательные файлы\textbackslash{}Актуальность\_адапт\_РНМ\textbackslash{}Ново-Коньковский 3-4СШ»)}}
    \label{fig:ch1_img8}
\end{figure}

Иногда же этот процесс приводит к ещё более серьёзным задержкам срабатывания, которые сложно оценить, так как длительность предаварийного режима осциллограммы на многих объектах 200-500~мс. Вот \commentnote{DISS-005-C0022}{Алексей Евдаков [2]}{t211131104284\_o0203}{пример }одной такой осциллограммы, где перед срабатывание частота была уже 26~Гц. На рисунке 4 вы можете видеть скриншоты частоты сигнала в начале осциллограммы и в момент итогового переключения.

\begin{figure}[htbp]
    \centering
    \includegraphics[width=0.85\textwidth]{image9.png}
    \caption[Частота сигнала на начало записи осциллограммы]{Частота сигнала на начало записи \commentnote{DISS-005-C0023}{Евдаков Алексей}{Улучшить картинку, здесь отмечены главные идеи}{осциллограммы}\newline{\small(Файл находится в «\_вспомогательные файлы\textbackslash{}Актуальность\_адапт\_РНМ»)}}
    \label{fig:ch1_img9}
\end{figure}

Из этой осциллограммы можно сделать, что иногда РНМ блокирует работу даже при отсутствии питания и почти полноценном останове двигательной нагрузки. И в итоге БАВР не может отработать, что можно считать ложным функционирование устройства.

При этом, «нагрузка» здесь остаётся симметричной и угол между фазным напряжением и током равен 90 градусов, даже с учётом того, что частота далека от 50Гц и измерения с помощью Фурье преобразований имеет большую погрешность. То есть РНМ остаётся в зоне потребления мощности (Срабатывание РНМ и соответственно блокировка). И поведение локально рассматриваемого органа с точки зрения заданных характеристик верное (рисунок 5).

\begin{figure}[htbp]
    \centering
    \begin{minipage}{0.48\textwidth}
        \centering
        \includegraphics[width=\textwidth]{image10.png}
        \centerline{\small а)}
    \end{minipage}\hfill
    \begin{minipage}{0.48\textwidth}
        \centering
        \includegraphics[width=\textwidth]{image11.png}
        \centerline{\small б)}
    \end{minipage}
    \caption[Осциллограмма: а) векторная диаграмма токов и напряжения; б) векторная диаграмма и характеристика РНМ]{Осциллограмма: а) \commentnote{DISS-005-C0024}{Евдаков Алексей}{Улучшить картинки}{векторная }диаграмма токов и напряжения; б) векторная диаграмма и характеристика РНМ\newline{\small(Файл находится в «\_вспомогательные файлы\textbackslash{}Актуальность\_адапт\_РНМ»)}}
    \label{fig:ch1_img10_11}
\end{figure}

\commentnote{DISS-005-C0025}{Алексей Евдаков [2]}{Улучшить текст, в рамках единого раздела.}{Скорее }всего в момент начала аварии была просадка амплитуды и изменение угла, но они оказались недостаточны, чтобы выйти за зону срабатывания. Это подтверждает важность и полезность адаптивной уставки (модернизации).

Описанная проблема возникает, когда уставка по модулю органа РНМ устанавливается на основе минимального порога чувствительности трансформаторов тока. Для её частичного решения, на ряде подстанций уставка РНМ подбиралась на основе минимально нагрузки на секциях шин. Однако это требует точной информации о характере нагрузки и её изменении. При этом, время работы БАВР всё равно остаётся не стабильным. Так как на большинстве подстанций нагрузка может сильно изменяться и в случае минимальной нагрузки -- время реакции РНМ будет мало, а в случае большой, может быть близко к указанным 20~мс даже с учётом подстройки и это в хороших случаях, а в плохих -- вообще приводить к несрабатыванию.

Подстройка порога срабатывания частично может быть решена за счёт применения различных групп уставок. Однако это требует ручного регулирования, и вводит дополнительное влияние человеческого фактора.

\noindent ДОБАВИТЬ НОРМАЛЬНОЕ \commentnote{DISS-005-C0026}{Алексей Евдаков [2]}{Просто накинул, чтобы не потерять\newline Беседы с Сашей от 28.08.2025, можно там поискать.}{\commentnote{DISS-005-C0027}{Алексей Евдаков [2]}{Терминал 931}{ОПИСАНИЕ}}

Так же при использовании нагрузки 0,4~кВ возникает огромное количество сложностей, к которому текущие устройства плохо приспособлены, а использование катушек Роговского добавляет к этому дополнительные сложности. Так как нагрузка может очень сильно колебаться, например на этом объекте от 90~А до 400~А, а угол от 0 до «-90» градусов, что ставит под сомнение саму идею использования. А ещё, оно может меняться в зависимости от того, какая нагрузка включена. Частотники прям очень сильно вредят всему этому.

Ниже примеры картинок (отображена амплитуда)

\begin{figure}[htbp]
    \centering
    \includegraphics[width=0.85\textwidth]{image12.png}
    \label{fig:ch1_img12}
\end{figure}

А вот там же другая нагрузка, тут разный период колебаний

\begin{figure}[htbp]
    \centering
    \includegraphics[width=0.85\textwidth]{image13.png}
    \label{fig:ch1_img13}
\end{figure}
